% ============================================================================
% 03_framework.tex — The Lever Space: A Conceptual Framework
% NEW writing per plan.md §3 (row 3). Number-free by design; all empirical
% magnitudes live in Sections 5–7. Taxonomy reconciliation per plan §5.7.
% ============================================================================

\section{The Lever Space: A Conceptual Framework}\label{sec:framework}

Fix the outcome rule: a strategy-proof direct mechanism in which truthful
reporting is a dominant strategy. Everything a designer can still change---the
extensive form through which reports are elicited, the language in which the
rules are explained, and the reasoning the interface invites the participant to
perform---constitutes the \emph{lever space} of this paper. The theory of simple
mechanisms tells us these residual choices are not innocuous. A second-price
sealed-bid (SPSB) auction and an ascending clock auction implement the same
outcome rule, yet only the clock is obviously strategy-proof
\citep{li2017obviously}; a mechanism and its description are formally distinct
objects, and descriptions can be engineered to expose strategyproofness
\citep{gonczarowski2023strategyproofness}; and the cognitive demands of a
mechanism decompose into distinct axes---contingent reasoning, forward
planning, and belief formation---that different designs tax differently
\citep{li2024designing}. The SPSB auction makes this lever space
fully operable. It is strategy-proof, and \emph{strategically simple} in the sense
of \citet{borgers2019strategically}---the dominant strategy is independent of
beliefs entirely---yet its direct form is not obviously strategy-proof, and it
admits an OSP implementation: the ascending clock
\citep{li2017obviously}. That wedge between strategy-proofness on paper and
obviousness in play is the identifying variation the levers act on. Auctions are
the paper's domain, and every lever family is evaluated there in the main text
(Section~\ref{sec:ranking}). One lever---the description of a fixed
mechanism---is additionally probed in a supplementary school-choice matching
experiment (student-proposing deferred acceptance under Ergin-acyclic
priorities \citep{ergin2002efficient,ashlagi2018stable}), reported in
Appendix~\ref{app:da}, where the menu description can be unbundled into its
invariance content and its mechanics; we draw on it in the main text only where
it sharpens the descriptions result.

We organize the levers into three main families---\textbf{A} (extensive form),
\textbf{B} (descriptions), and \textbf{C} (cognitive scaffolds)---plus an
auxiliary family \textbf{D} (preference framings) reported in
Appendix~\ref{app:prospect}. The taxonomy, and the predicted sign of every
lever in it, is derived ex ante from published theory rather than from
exploratory prompt search; roughly half the taxonomy is predicted to be inert
or to \emph{backfire}, and those predictions are part of what the experiments
test. Table~\ref{tab:framework-levers} summarizes the space;
Section~\ref{sec:design} gives the operational treatments and
Appendix~\ref{app:prompts} the verbatim prompt texts.

\subsection{Three Families of Levers}\label{sec:framework-families}

\paragraph{Family A: Extensive-form implementation.}
The strongest lever changes the game itself while holding the outcome rule
fixed. A mechanism is obviously strategy-proof (OSP) if, at every information
set where a deviation first diverges from the dominant strategy, the
\emph{worst} outcome from following the dominant strategy is at least as good
as the \emph{best} outcome from the deviation \citep{li2017obviously}.
Verifying dominance in an OSP mechanism therefore requires no hypothetical,
case-by-case reasoning about unobserved opponent moves: the comparison can be
made on-path, one information set at a time. The ascending clock auction is the
canonical example---it is OSP where SPSB is not, and it is additionally
\emph{one-step simple} in the sense of \citet{pycia2023theory}: comparing
``stay one more tick'' against ``exit now'' suffices at each price, with no
backward induction over the price trajectory. Our Family~A lever is the
ascending clock, in open AC and closed AC-B variants (Section~\ref{sec:design}).

\paragraph{Family B: Mechanism descriptions.}
The second family holds the extensive form fixed and changes only the words.
\citet{gonczarowski2023strategyproofness} formalize
\emph{strategyproofness-exposing} descriptions: honest accounts of a mechanism
from which its incentive properties are transparent. We distinguish three
sub-families, because they turn out to behave very differently.
\textbf{(B1) Menu restatement} re-describes the outcome function in menu form:
the opponent bids fix a menu (a price at which you may win), and your report
merely selects from it. This is the description
\citet{gonczarowski2023strategyproofness} propose for second-price auctions; in
incentivized human experiments, a menu-based explanation of DA improved
measured \emph{understanding} of strategyproofness while average
\emph{behavioral} effects were modest \citep{gonczarowski2024describing}.
In the matching domain we deploy the menu description in two variants---one
that includes the invariance statement (your report cannot change the menu you
face, only which item you receive from it), and one that restates only the
mechanics by which the menu is computed---so that the menu \emph{format} can
be separated from its invariance \emph{content}. The distinction is drawn ex
ante by \citet{gonczarowski2023strategyproofness}, who separate descriptions
engineered so that incentive properties are transparent from the text from
alternative, equally correct accounts of the same outcome function.
\textbf{(B2) Safety-exposing descriptions} state, in one or two sentences, the
structural invariant that makes truthful play safe. In the auction, this is the
pivotal logic at the heart of \citet{vickrey1961counterspeculation}'s original
argument: your bid determines \emph{whether} you win, never \emph{what} you pay
(\emph{Payoff Safety}). In DA, it is the property that historically
distinguished DA from the Boston mechanism
\citep{abdulkadirouglu2003school}: rejections only redirect, never
eliminate---your ranking determines the order in which you are considered, not
whether you are considered (\emph{Rejection Safety}). In the matching domain
we additionally test a plain textbook statement of strategy-proofness itself,
the bluntest member of this sub-family (Appendix~\ref{app:da}). \textbf{(B3)
Clock-framing} describes the sealed bid as an automatic exit threshold in a
hypothetical ascending clock---``your bid is the price at which you would stop
bidding''---following \citet{breitmoser2022obviousness}, who show that such
OSP-style \emph{descriptions} of SPSB increase truthful bidding in humans
without any change to the mechanism. We emphasize the accounting: clock-framing
is a description of a sealed-bid game and belongs to Family~B; the actual
ascending clock is a different extensive form and belongs to Family~A. The two
are never pooled in our analysis.

\paragraph{Family C: Cognitive scaffolds.}
The third family holds both the mechanism and its rule description fixed and
instead prompts the agent to \emph{reason} in a particular way---scaffolding
without revealing the dominant strategy. We index the scaffolds by the three
cognitive axes of \citet{li2024designing}. \textbf{(C1) Contingent reasoning}:
the core demand of non-OSP strategy-proof mechanisms is case-by-case reasoning
over unobserved opponent moves---``calculating payoffs for each profile of
opponent bids'' \citep{li2024designing}. Our \emph{Payoff Tree} scaffold
performs this construction for the agent, laying out the explicit tree of
payoff contingencies (win and pay the second-highest bid; lose and pay
nothing). It says what \emph{happens} under each action, never
what to \emph{do}. \textbf{(C2) Forward planning}: \citet{pycia2023theory}
formalize limited foresight over an agent's own future decision points. A
one-shot sealed bid has no such axis---there are no future own moves to plan
over---so forward planning has no place to bind in our auction domain and does
not enter the ranking. \textbf{(C3) Belief formation}: since the mechanism is
strategy-proof, optimal play is independent of beliefs about opponents---this
is precisely strategic simplicity \citep{borgers2019strategically}. First-order (``what do
you expect the others to do?'') and second-order (``what do the others think
\emph{you} will do?'') belief scaffolds are therefore irrelevant to a
participant who has understood dominance. Their inclusion is deliberately
diagnostic, as we formalize in H3 below.

\paragraph{Family D: Preference framings (appendix).}
Finally, we test framings that manipulate induced preferences rather than
reasoning: risk-attitude personas and prospect-theoretic loss/endowment frames
\citep{kahneman1979prospect, dreyfuss2022expectations}. Under a dominant
strategy, these are payoff-irrelevant to optimal play---bidding one's value is
dominant for any monotone utility---so they should be inert for a rational
agent. They are not levers a designer would pull to \emph{improve} play, but
they complete the bottom rungs of the ranking and calibrate how much damage an
adversarial or careless framing can do; results are in
Appendix~\ref{app:prospect}.

\begin{table}[t]
\centering
\footnotesize
\begin{tabular}{@{}p{0.26\linewidth}p{0.12\linewidth}p{0.24\linewidth}p{0.28\linewidth}@{}}
\toprule
Lever (family) & Domain & Theoretical anchor & Ex-ante prediction \\
\midrule
Ascending clock, AC/AC-B (A) & Auction & \citet{li2017obviously}; \citet{pycia2023theory} & Largest improvement \\
\addlinespace
Menu restatement (B1) & Auction, DA\textsuperscript{a} & \citet{gonczarowski2023strategyproofness} & Weak/none (human behavioral null) \\
Payoff / Rejection Safety (B2) & Auction & \citet{gonczarowski2023strategyproofness}; \citet{vickrey1961counterspeculation} & Improvement, approaching OSP \\
Clock-framing (B3) & Auction & \citet{breitmoser2022obviousness} & Improvement, approaching OSP \\
\addlinespace
Payoff Tree (C1) & Auction & \citet{li2024designing} & Improvement (axis binds) \\
Beliefs, 1st/2nd order (C3) & Auction & \citet{borgers2019strategically} & Inert if dominance understood; any effect diagnoses mis-strategizing \\
\addlinespace
Risk / loss framings (D) & Auction & \citet{kahneman1979prospect} & Payoff-irrelevant; inert for a rational agent (Appendix~\ref{app:prospect}) \\
\bottomrule
\end{tabular}
\caption{The lever space. Every lever holds the outcome rule fixed; families differ in what else they hold fixed (Family A changes the extensive form; B changes only the description; C changes only the invited reasoning; D changes only the induced preferences). Predictions are stated ex ante from the cited theory. All lever cells are reported in the auction main text (Section~\ref{sec:ranking}); the menu description is additionally probed in a supplementary matching experiment (Appendix~\ref{app:da}). Verbatim prompt texts and the old-to-new label mapping appear in Appendix~\ref{app:prompts}. \textsuperscript{a}\,In the matching experiment the menu description is deployed in two variants---invariance-property and mechanics-only---to separate the menu format from its invariance content (Appendix~\ref{app:da}).}
\label{tab:framework-levers}
\end{table}

\subsection{Hypotheses}\label{sec:framework-hypotheses}

The theory above yields three ordered hypotheses, which we state before
presenting any results. They pre-commit the paper's predictions and, because
several predict null or negative effects, they discipline the exercise against
the objection that prompt variations were searched until something worked.

\paragraph{H1 (Extensive form dominates).}
Family~A levers produce the largest and most robust improvements in truthful
play. They are the only levers that change the game the agent actually faces:
under OSP, verifying dominance no longer requires contingent reasoning at all
\citep{li2017obviously}, and the clock is additionally one-step simple
\citep{pycia2023theory}. Every other family asks the agent to do the same hard
verification with better hints; Family~A removes the verification problem.

\paragraph{H2 (Descriptions substitute when---and only when---they expose the
invariant).}
Within Family~B, descriptions that state the invariant making truthful play
safe---the safety descriptions (B2), clock-framing (B3), and the
invariance-property variant of the menu---approach the OSP benchmark
\emph{without} any change to the extensive form, while descriptions that
merely restate the outcome function in different words---the plain menu
restatement and its mechanics-only variant---do not, and may even backfire.
The operative distinction is thus \emph{invariance-exposing} versus
\emph{mechanically-restating} text, a line drawn ex ante by
\citet{gonczarowski2023strategyproofness} between descriptions engineered so
that incentive properties are transparent and alternative, equally correct
accounts of the same outcome function. The logic: if the binding constraint is
not \emph{computing} the dominant strategy but \emph{seeing that truthful play
is safe}, then supplying the safety certificate verbally should substitute for
supplying it structurally---while supplying more words without the certificate
should not. The mechanical restatements are the built-in placebo for this
claim: they change wording, length, and salience while adding no invariant.
The human evidence already points this way: menu descriptions improve measured
comprehension with modest behavioral effects \citep{gonczarowski2024describing},
whereas OSP-style clock descriptions shift human bids toward truthfulness
\citep{breitmoser2022obviousness}. A confirmation of H2 would be the paper's
most useful design result, since descriptions are the cheapest lever a
practitioner can pull.

\paragraph{H3 (Scaffolds are axis-specific; belief scaffolds are diagnostic).}
A Family~C scaffold helps only on the axis where the mechanism's cognitive
demand actually binds. Concretely: (i) C1 (contingent reasoning) should help,
because case-by-case reasoning over opponent profiles is the core
demand of the non-OSP sealed-bid format \citep{li2024designing}; (ii) C2 (forward
planning) has no axis to relax in a one-shot sealed bid, so it does not enter
the auction ranking; (iii) C3 (beliefs) should be exactly inert, because
dominant strategies are belief-free \citep{borgers2019strategically}. The
diagnostic content of H3 is in its failure modes. Any measured effect of a
belief scaffold---in either direction---is evidence that the agent is not
playing the dominant strategy but treating the interaction as a Bayesian game,
forming beliefs about competitors and best-responding to them. Belief scaffolds
that activate strategic simulation without resolving it hand the agent material
to strategize with; if the agent has not internalized strategy-proofness, we
predict such scaffolds \emph{backfire}, and most severely for the weakest
reasoners. Backfiring is thus not an embarrassment for
the taxonomy but its theoretically expected signature: a designer's lever space
contains anti-levers, and knowing which is which is part of the ranking this
paper delivers (Section~\ref{sec:ranking}).

\paragraph{A note on taxonomy provenance.}
The operational labels in our experiment repository predate the final taxonomy,
and three assignments were corrected during its construction (the full
old-to-new mapping is in Appendix~\ref{app:prompts},
Table~\ref{tab:intervention-inventory}). First, the \emph{Payoff Tree} scaffold
was originally filed under forward planning (repo label
\texttt{axis2\_forward\_tree}); it is classified here as contingent reasoning
(C1), because the tree enumerates payoff \emph{contingencies}---win versus
lose, and what each implies---not future own moves, and a one-shot sealed bid
has no forward-planning axis for it to occupy. Second, \emph{Rejection Safety}
was drafted within a DA ``monotonicity'' scaffold family
(\texttt{axis2\_da\_monotonic\_safety}); it scaffolds no reasoning at
all---it states a property of the mechanism---and is therefore a
safety-exposing description (B2). Third, one repository template name
(\texttt{intervention\_proxy\_breitmoser}) was historically shared between the
sealed-bid clock-framing \emph{description} and a true iterative clock; in this
paper the two are strictly separated as B3 and Family~A respectively, with
distinct prompt identifiers (\texttt{desc\_clock\_framing} vs.\
\texttt{osp\_clock\_iterative}) throughout. In the harmonized cell dataset
(Section~\ref{sec:design}) the two conditions appear as separate cells for all
four models---the sealed-bid clock-framing description and the closed
ascending clock---so rungs A and B3 are now distinguished empirically as well
as by design.

With the lever space fixed and its predictions on record, the next section
turns the taxonomy into an experimental design: environments, treatment
implementations, models, metrics, and the human benchmarks against which the
testbed is validated.
